\documentclass[11pt]{article}
\usepackage[a4paper,margin=25mm]{geometry}
\usepackage{amsmath,amssymb,amsthm,mathtools}
\usepackage{graphicx}
\usepackage[hidelinks]{hyperref}
\usepackage{xcolor}
\usepackage{microtype}
\usepackage{enumitem}

\definecolor{accent}{RGB}{36,80,112}
\newtheorem{theorem}{Theorem}
\newtheorem{corollary}[theorem]{Corollary}
\newtheorem{remark}[theorem]{Remark}

\title{Canonical-Dual Sobolev Regularity:\\
\large A sharp shift obstruction and a positive low-regularity theorem}
\author{Automated mathematics run \texttt{fa\_banach\_001}, lane 8}
\date{17 August 2026}

\begin{document}
\maketitle

\begin{center}
\fcolorbox{accent}{accent!6}{\parbox{0.91\textwidth}{
\textbf{Status.} Candidate substantial partial result, likely valid, pending
human review. The concrete boundary-shearlet estimate at the intended
integer Sobolev regularities remains open.}}
\end{center}

\section{Source problem}

Grohs, Kutyniok, Ma, Petersen, and Raslan construct boundary shearlet frames
$\Phi=(\varphi_n)$ for $L^2(\Omega)$. Writing $j_n$ for the scale, $C$ and
$D$ for primal analysis and synthesis, and $S=DC$ for the $L^2$ frame
operator, their desired Gelfand-frame estimates at regularity $s>0$ are
\begin{align*}
 \left\|D c\right\|_{H^s(\Omega)}
   &\lesssim \left\|(2^{j_ns}c_n)_n\right\|_2, \tag{GFA1}\\
 \left\|(2^{j_ns}\langle f,S^{-1}\varphi_n\rangle)_n\right\|_2
   &\lesssim \|f\|_{H^s(\Omega)}. \tag{GFA2}
\end{align*}
They prove a weighted primal-analysis characterization of $H^s(\Omega)$ and
derive (GFA1) under localization assumptions. They explicitly leave (GFA2)
without an analytic proof and provide numerical evidence instead
\cite[pp.~5, 18, and \S6.3]{GKMPMR}.

\begin{figure}[ht]
  \centering
  \includegraphics[width=0.97\textwidth]{figures/open_problem_crop.png}
  \caption{The source explains that (GFA2) is inaccessible because it concerns
  canonical-dual coefficients, and records only a numerical analysis
  \cite[p.~18]{GKMPMR}.}
\end{figure}

The first theorem below proves that the two available mapping properties do
not imply (GFA2), even under very strong qualitative localization. The second
theorem gives a positive analytic interval around $s=0$.

\section{Proof intuition}

The obstruction is directional. A unilateral shift to a higher scale costs
a factor $2^s$ in the weighted norm, while the backward shift gains the same
factor. Choose a Riesz-basis synthesis operator $D=I-aR$. Its adjoint
$D^*=I-aR^*$ controls the primal coefficients extremely well because it moves
toward lower scales. Canonical-dual analysis, however, is $D^{-1}$ and sums
an infinite forward-shift tail. That tail belongs to the weighted space
exactly below the threshold $a2^s=1$.

For a general frame the frame operator $S$ is nevertheless self-adjoint and
invertible on the central $L^2$ space. Endpoint analysis and synthesis make
$S$ bounded at positive regularity, and self-adjointness supplies the negative
endpoint by duality. Sneiberg's interpolation theorem then preserves
invertibility on some open interval around zero. This proves (GFA2) there,
although it gives no reason for invertibility to persist all the way to a
prescribed integer exponent.

\section{An exact phase transition}

Let $\mathcal H_0=\ell^2(\mathbb N_0)$ with standard basis $(e_n)$, let
$R e_n=e_{n+1}$, and for $r\geq0$ put
\[
 \mathcal H_r=\left\{x:\|x\|_r^2=
 \sum_{n\geq0}2^{2rn}|x_n|^2<\infty\right\}.
\]
Use the same weighted space for coefficients and denote it by $\ell^2_r$.

\begin{theorem}[Bidiagonal Riesz-basis classification]\label{thm:shift}
Fix $0<a<1$, set $D=I-aR$, and define $\varphi_n=De_n$. Then
$(\varphi_n)$ is a Riesz basis of $\mathcal H_0$. For every $r\geq0$:
\begin{enumerate}[label=\textup{(\roman*)},leftmargin=2.1em]
\item $D:\ell^2_r\to\mathcal H_r$ is bounded, so \emph{(GFA1)} holds;
\item its primal coefficients give an exact scale characterization,
\[
 \|(\langle f,\varphi_n\rangle)_n\|_{\ell^2_r}
 =\|D^*f\|_r\asymp\|f\|_r;
\]
\item canonical-dual analysis is bounded $\mathcal H_r\to\ell^2_r$ if and
only if
\[
 a2^r<1.
\]
In particular, (GFA1) and the primal coefficient characterization do not imply
(GFA2), even for a Riesz basis with tridiagonal Gramian and frame operator.
\end{enumerate}
\end{theorem}

\begin{proof}
Because $\|aR\|_{\mathcal H_0}=a<1$, the Neumann series makes $D$ invertible
on $\mathcal H_0$. Thus $(De_n)$ is a Riesz basis, with ordinary frame bounds
between $(1-a)^2$ and $(1+a)^2$.

Put $q=2^r$. Directly from the weighted norm,
\[
 \|R\|_{\mathcal H_r}=q,
 \qquad \|R^*\|_{\mathcal H_r}=q^{-1}.
\]
Consequently $\|Dc\|_r\leq(1+aq)\|c\|_r$, proving (i). Primal
analysis is $D^*=I-aR^*$, and
\[
 (D^*)^{-1}=\sum_{k\geq0}a^k(R^*)^k
\]
converges in $\mathcal B(\mathcal H_r)$ because $a/q<1$. Equivalently,
\[
 (1-a/q)\|f\|_r\leq\|D^*f\|_r
 \leq(1+a/q)\|f\|_r,
\]
which proves (ii).

The canonical dual basis is $((D^{-1})^*e_n)_n$, so its analysis map is
$D^{-1}$. If $aq<1$, the Neumann series
$D^{-1}=\sum_{k\geq0}a^kR^k$ converges on $\mathcal H_r$. Conversely,
\[
 D^{-1}e_0=\sum_{k\geq0}a^k e_k,
 \qquad
 \|D^{-1}e_0\|_r^2=\sum_{k\geq0}(aq)^{2k},
\]
which diverges when $aq\geq1$. This proves (iii). Since $D$ is bidiagonal,
$D^*D$ and $DD^*$ are tridiagonal.
\end{proof}

\begin{corollary}[Near-Parseval failure at fixed smoothness]
Fix $r>0$ and $\eta>0$. There is a scale-labelled Riesz basis with ordinary
$L^2$ frame bounds in $[(1-\eta)^2,(1+\eta)^2]$ for which (GFA1) and the
primal characterization hold at $r$, but (GFA2) fails.
\end{corollary}

\begin{proof}
Label $e_n$ by scale $j_n=Ln$. The weight ratio is then $q=2^{Lr}$.
Choose $L$ so large that $a=2^{-Lr}<\min\{\eta,1\}$. Theorem~\ref{thm:shift}
applies with $aq=1$.
\end{proof}

\section{A positive interval for every endpoint-stable frame}

Let $(\mathcal H_t)_{-\sigma\leq t\leq\sigma}$ and
$(\ell^2_t)_{-\sigma\leq t\leq\sigma}$ be exact complex interpolation
Hilbert scales, with $\mathcal H_{-t}=\mathcal H_t'$ under the central
$\mathcal H_0$ pairing. Standard Sobolev spaces and dyadically weighted
$\ell^2$ spaces have these properties.

\begin{theorem}[Low-regularity Gelfand interval]\label{thm:interval}
Let $\Phi$ be a frame for $\mathcal H_0$, with analysis $C$, synthesis $D$,
and frame operator $S=DC$. Suppose for some $\sigma>0$ that
\[
 C:\mathcal H_\sigma\to\ell^2_\sigma,
 \qquad D:\ell^2_\sigma\to\mathcal H_\sigma
\]
are bounded. Then there exists $\varepsilon\in(0,\sigma]$ such that, for
every $0\leq r<\varepsilon$, both (GFA1) and (GFA2) hold at exponent $r$.
Equivalently, $\Phi$ is a Gelfand frame throughout a nonzero interval of
regularities adjacent to $L^2$.
\end{theorem}

\begin{proof}
The hypotheses imply that $S=DC$ is bounded on $\mathcal H_\sigma$. Since
$S$ is self-adjoint for the $\mathcal H_0$ pairing, duality extends it
boundedly to $\mathcal H_{-\sigma}$. It is boundedly invertible on
$\mathcal H_0$ because $\Phi$ is an $\mathcal H_0$ frame.

Sneiberg's stability theorem for operators on complex interpolation scales
now gives $\varepsilon>0$ such that
$S:\mathcal H_r\to\mathcal H_r$ is boundedly invertible whenever
$|r|<\varepsilon$ \cite{Sneiberg}. Interpolation between the central and
positive endpoints gives bounded maps
$C:\mathcal H_r\to\ell^2_r$ and
$D:\ell^2_r\to\mathcal H_r$ for $0\leq r\leq\sigma$. Hence GFA1 holds,
and canonical-dual analysis satisfies
\[
 C_{\Phi^d}f=C S^{-1}f,
 \qquad
 \|C_{\Phi^d}f\|_{\ell^2_r}
 \leq \|C\|\,\|S^{-1}\|\,\|f\|_{\mathcal H_r}
\]
for $0\leq r<\varepsilon$. This is GFA2.
\end{proof}

\begin{remark}
Theorem~\ref{thm:interval} applies to a boundary-shearlet system whenever the
source's endpoint coefficient characterization and localized synthesis bound
are available at some $\sigma>0$. It does not quantify $\varepsilon$ from
the published estimates and therefore need not reach $r=1$.
\end{remark}

\section{Consequences, limitation, and verification}

Theorem~\ref{thm:shift} identifies the exact logical gap in the source's
program. Boundedness of the Gramian on the weighted space, finite-band
localization, and even near-tight $L^2$ frame bounds do not ensure that the
canonical inverse respects the weight. A full boundary-shearlet proof needs
a quantitative inverse-closed localization algebra whose decay dominates
$2^{sj}$; this is stronger than the componentwise weighted boundedness and
the particular cross-decay estimate proved in the source. General
inverse-closed localization results explain the appropriate route
\cite{GK09,FSS19}, but verifying their hypotheses for the full hybrid index
geometry was not achieved here.

All computations in Theorem~\ref{thm:shift} are exact geometric-series and
operator-norm identities. The low-regularity theorem uses only frame
operator self-adjointness, Hilbert-scale duality, interpolation of $C,D$, and
Sneiberg stability. Eight focused routes were audited: the exact shift
classification, near-Parseval strengthening, and low-regularity theorem
succeeded; tight-frame reduction, direct implication, offset-Neumann,
inverse-closed hybrid localization, and later-literature closure did not
settle the intended integer-regularity problem.

The run's four lightweight indexes were searched by arXiv id, title,
'GFA2', 'boundary shearlet', 'canonical dual', and 'Gelfand frame'.
A bounded arXiv/web search through 17 August 2026 found the source, its 2017
follow-up \cite{PR17}, and general localized-frame literature, but no analytic
resolution for the concrete boundary-shearlet system or exact match for the
shift phase transition. This is bounded evidence, not a priority claim.

\paragraph{Human-review recommendation.}
Promote for expert review as a substantial partial result. Check the shift
orientation and canonical-dual identity in Theorem~\ref{thm:shift}, and check
the interpolation-couple hypotheses and use of Sneiberg's theorem in
Theorem~\ref{thm:interval}. Do not promote this packet as a full answer to
the boundary-shearlet question without the missing inverse-closed localization
estimate.

\begin{thebibliography}{9}
\small
\bibitem{GKMPMR}
P. Grohs, G. Kutyniok, J. Ma, P. Petersen, and M. Raslan,
\emph{Anisotropic Multiscale Systems on Bounded Domains},
arXiv:1510.04538 (2015; revised 2017), Adv. Comput. Math. 46 (2020), Art.~39.

\bibitem{PR17}
P. Petersen and M. Raslan,
\emph{Approximation properties of hybrid shearlet--wavelet frames for Sobolev
spaces}, arXiv:1712.01047 (2017).

\bibitem{Sneiberg}
I. Ya. Sneiberg,
\emph{Spectral properties of linear operators in interpolation families of
Banach spaces}, Mat. Issled. 9 (1974), no.~2, 214--229.

\bibitem{GK09}
K. Gr\"ochenig and A. Klotz,
\emph{Noncommutative approximation: inverse-closed subalgebras and
off-diagonal decay of matrices}, Constr. Approx. 32 (2010), 429--466;
arXiv:0904.0386.

\bibitem{FSS19}
Q. Fang, C. E. Shin, and Q. Sun,
\emph{Polynomial control on weighted stability bounds and inversion norms of
localized matrices on simple graphs}, arXiv:1909.08409 (2019).
\end{thebibliography}

\end{document}

